\documentclass[a4paper,11pt]{article}
\usepackage[utf8]{inputenc}
\usepackage[T1]{fontenc}
\usepackage{lmodern}
\usepackage{amsmath,amssymb,amsthm}
\usepackage{mathtools}
\usepackage{booktabs}
\usepackage{longtable}
\usepackage{array}
\usepackage{hyperref}
\usepackage{graphicx}
\usepackage{float}
\usepackage{geometry}
\geometry{margin=25mm}

\hypersetup{
  colorlinks=true,
  linkcolor=blue,
  urlcolor=blue,
  citecolor=blue
}

\theoremstyle{definition}
\newtheorem{theorem}{Theorem}[section]
\newtheorem{proposition}[theorem]{Proposition}
\newtheorem{lemma}[theorem]{Lemma}
\newtheorem{corollary}[theorem]{Corollary}
\newtheorem{definition}[theorem]{Definition}
\newtheorem{remark}[theorem]{Remark}
\newtheorem{observation}[theorem]{Observation}

\setlength{\parindent}{1em}
\setlength{\parskip}{0.5em}

\providecommand{\tightlist}{\setlength{\itemsep}{0pt}\setlength{\parskip}{0pt}}
\begin{document}

\section{Subjective Spaces from Zero to Four Dimensions}\label{subjective-spaces-from-zero-to-four-dimensions}

\begin{center}\rule{0.5\linewidth}{0.5pt}\end{center}

\textbf{Author}: Noriaki Kihara (WF System Co., Ltd.)

\textbf{Date}: April 2026

\textbf{DOI:} \href{https://doi.org/10.5281/zenodo.19533292}{10.5281/zenodo.19533292}

\begin{center}\rule{0.5\linewidth}{0.5pt}\end{center}

\subsection{Abstract}\label{abstract}

Based on the central projection framework and the interpretive structure of dimensional extension established in prior papers {[}1{]}, {[}2{]}, {[}3{]}, {[}4{]}, {[}5{]}, this paper describes the geometric properties that emerge at each dimension of subjective space from zero to four dimensions, and their correspondence to the structures regarded as standard in each case.

\textbf{Clarification}: As established in prior papers {[}4{]} and {[}5{]}, dimensional extension within the central projection framework is in principle executable without limit. How far dimensions are added, and how the stopping condition for extension is to be interpreted, are matters of interpretation and lie outside the scope of this paper. For convenience of exposition, this paper treats subjective spaces from zero to four dimensions; however, the number four is merely a notational convenience, and this paper does not claim any geometric reason why dimensional extension should stop at four dimensions.

This paper proves nothing new. It only describes how the geometric properties regarded as standard in a given dimension correspond to the subjective space of the central projection framework at that dimension. Physical interpretation lies outside the scope of this paper and is left to the reader.

\begin{center}\rule{0.5\linewidth}{0.5pt}\end{center}

\subsection{1. Introduction}\label{introduction}

In the central projection framework introduced in prior paper {[}1{]}, the subjective space \(\Pi_R\) is defined as a space governed by the curvature radius \(R\). Paper {[}2{]} demonstrated the equivalence among the axes of subjective space, paper {[}3{]} demonstrated the relativity of observation across multiple subjective spaces, paper {[}4{]} demonstrated the existence of limits of the curvature radius and nested structures, and paper {[}5{]} demonstrated the interpretive constraints on dimensional extension.

Building on these results, this paper examines the dimensions of subjective space in sequence: \(0, 1, 2, 3, 4\). For each dimension, it describes the geometric properties that newly emerge at that dimension, the quantities that become cognizable, and observations regarding the interpretation of the time axis.

To reiterate, this paper proves nothing new about these correspondences. The description at each dimension is simply the application of the geometric properties regarded as standard for that dimension to the subjective space. The intent of this paper is limited to organizing the correspondence between the central projection framework and established theory.

\begin{center}\rule{0.5\linewidth}{0.5pt}\end{center}

\subsection{2. Zero-Dimensional Subjective Space}\label{zero-dimensional-subjective-space}

The zero-dimensional subjective space is a world that distinguishes only ``presence / absence.''

\textbf{Internal structure}: None.

\textbf{Cognizable quantities}: Geodesic deviation does not exist. There is also no means of recognizing from the inside whether one is governed by a curvature radius \(R\).

\textbf{Interpretation of the time axis}: Not applicable (no axes exist).

\textbf{Possibility of dimensional extension}: In accordance with prior paper {[}5{]}, dimensions may be freely added.

\begin{center}\rule{0.5\linewidth}{0.5pt}\end{center}

\subsection{3. One-Dimensional Subjective Space}\label{one-dimensional-subjective-space}

The one-dimensional subjective space is a world consisting solely of non-negative real numbers.

\textbf{Internal structure}: A single axis. However, neither a reference point nor a criterion for determining direction exists. It is a pure set of non-negative real numbers. The concept of length exists, but the concept of direction does not.

\textbf{Cognizable quantities}: Geodesic deviation is not yet cognizable (in one dimension, there is no other axis with respect to which deviation can be defined).

\textbf{Interpretation of the time axis}: Not applicable (there is only one axis, and no basis for comparison exists).

\textbf{Possibility of dimensional extension}: In accordance with prior paper {[}5{]}, dimensions may be freely added.

\begin{center}\rule{0.5\linewidth}{0.5pt}\end{center}

\subsection{4. Two-Dimensional Subjective Space}\label{two-dimensional-subjective-space}

In the two-dimensional subjective space, area and angle newly emerge.

\textbf{Internal structure}: Two axes. A surface is defined, and the surface acquires an orientation (front and back). Through this orientation of the surface, \textbf{direction} is conferred upon the one-dimensional length. The Pythagorean theorem holds.

\textbf{Cognizable quantities}: Area, angle, and geodesic deviation become cognizable. Geodesic deviation becomes definable for the first time in two dimensions, and the corresponding proposition in paper {[}2{]} acquires meaning.

\textbf{Interpretation of the time axis}: The two axes are equivalent by Theorem 3.1 of paper {[}2{]}, and it is not possible to determine which axis is to be interpreted as the time axis.

\textbf{Possibility of dimensional extension}: In accordance with prior paper {[}5{]}, dimensions may be freely added.

\begin{center}\rule{0.5\linewidth}{0.5pt}\end{center}

\subsection{5. Three-Dimensional Subjective Space}\label{three-dimensional-subjective-space}

In the three-dimensional subjective space, volume newly emerges.

\textbf{Internal structure}: Three axes. Volume is defined, and the area acquires an orientation (front and back).

\textbf{Cognizable quantities}: Volume, area, angle, and geodesic deviation are cognizable. Geodesic deviation has a richer structure in three dimensions.

\textbf{Interpretation of the time axis}: If one of the three axes is interpreted as the time axis, the remaining two are recognized as spatial axes. Which axis is interpreted as the time axis is left to the observer's choice by Theorem 3.1 of paper {[}2{]}.

\textbf{Possibility of dimensional extension}: In accordance with prior paper {[}5{]}, dimensions may be freely added.

\begin{center}\rule{0.5\linewidth}{0.5pt}\end{center}

\subsection{6. Four-Dimensional Subjective Space}\label{four-dimensional-subjective-space}

In the four-dimensional subjective space, a four-dimensional hypervolume newly emerges.

\textbf{Internal structure}: Four axes. A four-dimensional hypervolume is defined. The three-dimensional volume acquires an orientation (front and back).

\textbf{Cognizable quantities}: Four-dimensional hypervolume, three-dimensional volume, area, angle, and geodesic deviation are cognizable.

\textbf{Interpretation of the time axis}: If one of the four axes is interpreted as the time axis, the remaining three are recognized as spatial axes. On the other hand, if two axes are interpreted as time axes, they cannot be geometrically distinguished from the remaining two spatial axes (by virtue of the axis equivalence of Theorem 3.1 in paper {[}2{]}). That is, the number of axes interpreted as time axes is not uniquely determined within the interior of the four-dimensional subjective space.

\textbf{Possibility of dimensional extension}: In accordance with prior paper {[}5{]}, dimensions may be freely added.

\begin{center}\rule{0.5\linewidth}{0.5pt}\end{center}

\subsection{7. Discussion}\label{discussion}

This paper has described the correspondence between the geometric properties that newly emerge at each dimension of subjective space from zero to four dimensions and the properties regarded as standard at that dimension.

The descriptions in this paper are limited to the following.

\begin{enumerate}
\def\labelenumi{\arabic{enumi}.}
\tightlist
\item
  The quantities that become newly definable at each dimension (area, angle, volume, hypervolume, and the corresponding orientation structures).
\item
  The status of geodesic deviation that becomes cognizable at each dimension.
\item
  The interpretability of the time axis at each dimension.
\item
  The fact that dimensional extension is free at each stage.
\end{enumerate}

The properties described in this paper for each dimension represent correspondences to the geometric structures regarded as standard for that dimension, applied to the subjective space, and do not contain new proofs. The physical significance these correspondences may carry is not addressed in this paper. In particular, questions such as how the orientation structure that emerges at each dimension manifests in the observer's cognition, and by what criterion the interpretation of the time axis is realized, are problems outside the geometric framework and beyond the scope of this paper.

The answers to these questions are left to the reader's judgment.

\begin{center}\rule{0.5\linewidth}{0.5pt}\end{center}

\subsection{8. Conclusion}\label{conclusion}

The subjective space of the central projection framework corresponds naturally to the geometric properties regarded as standard at each dimension from zero to four dimensions. Dimensional extension is free at each stage, and the description up to four dimensions in this paper is purely a matter of expository convenience. The physical interpretation of the structures appearing at each dimension lies outside the scope of this paper and is left to the reader.

\begin{center}\rule{0.5\linewidth}{0.5pt}\end{center}

\subsection{References}\label{references}

{[}1{]} Noriaki Kihara, ``Geometric Formulation via Central Projection,'' Zenodo, DOI: 10.5281/zenodo.19427780 (2025).

{[}2{]} Noriaki Kihara, ``Symmetry in the Central Projection Framework,'' Zenodo, DOI: 10.5281/zenodo.19434932 (2025).

{[}3{]} Noriaki Kihara, ``Relativity of Observation in Multiple Subjective Spaces,'' Zenodo, DOI: 10.5281/zenodo.19435162 (2025).

{[}4{]} Noriaki Kihara, ``On the Limits of the Curvature Radius in Central Projection: The Cases \(R \to \infty\) and \(R \to 0\),'' Zenodo, DOI: 10.5281/zenodo.19526549 (2026).

{[}5{]} Noriaki Kihara, ``On Dimensional Extension of Background Space and Subjective Space,'' Zenodo, DOI: 10.5281/zenodo.19526913 (2026).

\begin{center}\rule{0.5\linewidth}{0.5pt}\end{center}

\end{document}
